
\begin{figure}[t]
%\setlength{\tabcolsep}{10pt}
\begin{center}
%\vspace{-1cm}
\includegraphics[width = .95\linewidth]{emnlp2023-latex/figures/dendrogram.pdf}
\end{center}
\caption{Performance of the transfer between cognitive and linguistic probes, obtained by training linguistic probes to detect specific linguistic features from language models embeddings, transferring those probes zero-shot to cognitive datasets, and vice versa. The probes trained on the behavioral and brain datasets are reported on the left, the linguistic probes on the bottom. The squares in the heatmap report the results of the transfer (Pearson correlation between predicted and actual responses), averaging across the two transfer directions (\textit{cognitive} $\leftrightarrows$ \textit{linguistic}). The bars on top of the heatmap indicate the average difference in transfer between behavioral and brain probes for each linguistic probe ($\Delta R$); the bars are colored in red if the transfer is significantly higher for behavioral probes, and blue if transfer is higher for brain probes (FDR-corrected). Behavioral probes obtain the strongest transfer performance with form-level probes, whereas brain probes transfer better with semantic probes.}
\label{fig:followup}

\end{figure}
